\documentclass[conference]{IEEEtran}
\IEEEoverridecommandlockouts

\usepackage{cite}
\usepackage{amsmath,amssymb,amsfonts}
\usepackage{algorithmic}
\usepackage{graphicx}
\usepackage{textcomp}
\usepackage{xcolor}
\usepackage{hyperref}
\usepackage{booktabs} % For professional tables

\begin{document}

\title{Unsupervised Statistical Learning for Kinematic Anomaly Detection in Autonomous Driving Scenarios}

\author{\IEEEauthorblockN{Pavel Stepanov}
\IEEEauthorblockA{\textit{Dept. of Computer Science} \\
\textit{University of Miami}\\
Coral Gables, USA \\
pas273@miami.edu}
\and
\IEEEauthorblockN{Ramses Loaces}
\IEEEauthorblockA{\textit{Dept. of Computer Science} \\
\textit{University of Miami}\\
Coral Gables, USA \\
rxl1238@miami.edu}
\and
\IEEEauthorblockN{Justin Cabrera}
\IEEEauthorblockA{\textit{Dept. of Computer Science} \\
\textit{University of Miami}\\
Coral Gables, USA \\
justin.cabrera2718@gmail.com}
}

\maketitle

\begin{abstract}
Autonomous vehicles must operate safely in environments characterized by stochastic human behavior. While nominal driving scenarios are well-represented in current datasets, identifying anomalous events---such as aggressive maneuvering, emergency braking, or near-miss interactions---remains a challenge due to the scarcity of labeled high-risk data. This paper presents an unsupervised statistical learning framework applied to the Waymo Open Motion Dataset (v1.3.1). We extract kinematic features from over 10,000 traffic agents and utilize Principal Component Analysis (PCA) for dimensionality reduction. We systematically compare parametric (Mahalanobis Distance) and non-parametric (One-Class Support Vector Machine) approaches for outlier detection. Experimental results demonstrate that statistically identified anomalies correlate strongly with ground-truth interaction flags provided by Waymo, validating the efficacy of unsupervised methods for filtering safety-critical scenarios.
\end{abstract}

\begin{IEEEkeywords}
Statistical Learning, Anomaly Detection, Waymo Open Dataset, Unsupervised Learning, PCA, One-Class SVM
\end{IEEEkeywords}

\section{Introduction}
The deployment of autonomous vehicles (AVs) requires rigorous validation against the "long tail" of driving distributions. Traditional supervised learning approaches are limited by the scarcity of labeled anomaly data. Consequently, unsupervised techniques are essential for mining large-scale datasets for safety-critical events. This study applies classical statistical learning methods to quantify kinematic deviance in urban traffic scenarios.

\section{Data Description}
We utilize the Waymo Open Motion Dataset (v1.3.1), specifically the validation split. The raw Protocol Buffer (TFRecord) data was processed to extract a structured tabular dataset consisting of $N \approx 10,000$ agents across 300 distinct scenarios.

\subsection{Feature Extraction}
For each agent, we extracted the following kinematic features at the prediction horizon $t=0$:
\begin{itemize}
    \item \textbf{Speed ($v$):} Magnitude of the velocity vector ($m/s$).
    \item \textbf{Acceleration ($a$):} Rate of change of speed ($m/s^2$).
    \item \textbf{Yaw Rate ($\dot{\psi}$):} Rate of change of heading ($rad/s$).
    \item \textbf{Lateral G-Force:} Calculated as $|v \times \dot{\psi}|$.
\end{itemize}

\section{Methodology}

\subsection{Preprocessing}
Data cleaning involved the removal of the autonomous vehicle (SDC) tracks to focus on human agent behavior. Missing values were handled via listwise deletion. Features were normalized using Z-score standardization to ensure scale invariance for distance-based metrics.

\subsection{Dimensionality Reduction}
We applied Principal Component Analysis (PCA) to reduce the dimensionality of the feature space while retaining maximum variance. The first two principal components were analyzed to identify clusters of nominal driving behavior.

\subsection{Anomaly Detection}
Two unsupervised methods were implemented:
\begin{enumerate}
    \item \textbf{Mahalanobis Distance:} A multivariate distance metric accounting for the covariance structure of the data. Observations exceeding the $\chi^2_{0.99}$ critical value were flagged as anomalies.
    \item \textbf{One-Class SVM:} A non-parametric approach utilizing a Radial Basis Function (RBF) kernel to define a decision boundary encapsulating the majority of the data distribution.
\end{enumerate}

\section{Results}
\subsection{Statistical Validation}
We validated our detected anomalies against the \texttt{IsInteractive} flag provided in the Waymo dataset. A contingency table analysis reveals a significant correlation between statistical outliers and interactive agents.

% Insert Figure Placeholder
% \begin{figure}[htbp]
% \centerline{\includegraphics[width=0.45\textwidth]{pca_biplot.png}}
% \caption{PCA Biplot of Kinematic Features.}
% \label{fig}
% \end{figure}

\section{Conclusion}
Our analysis confirms that unsupervised statistical learning is a viable methodology for identifying high-risk driving behaviors without the need for manual annotation. The Mahalanobis Distance metric proved particularly effective at isolating agents exhibiting physically extreme maneuvers.

\bibliographystyle{IEEEtran}
% \bibliography{references}

\end{document}